\begin{problem}[Sections of a sheaf are determined by their germs]
Let $\mathcal F$ be a sheaf and $s,t\in\mathcal F(U)$. Prove
\[
s=t\iff s_x=t_x\ \text{for all }x\in U.
\]
Show by counterexample that this statement can fail for presheaves that are not sheaves.
\end{problem}

\paragraph{Why this problem matters.}
This is one of the central local-to-global uses of stalks and the key reason stalks detect morphisms.

\paragraph{Useful ideas before starting.}
Equality of germs gives a neighborhood on which the sections agree. For the counterexample, use a presheaf with invisible global data.

\paragraph{Guided hints.}
For a counterexample, define a presheaf on a two-point discrete space with a nonzero group on the whole space and zero on each proper open.

\ifdefined\FullProblemDossiers
\begin{solution}
If $s=t$, the germs agree. Conversely, for every $x$ choose $V_x\subseteq U$ with $s|_{V_x}=t|_{V_x}$. These neighborhoods cover $U$, and sheaf locality gives $s=t$. For failure in a presheaf, let $X=\{p,q\}$ discrete, put $\mathcal P(X)=A\neq0$ and $\mathcal P(V)=0$ for every proper open $V$, with all proper restrictions zero. Distinct global elements have zero germs at both points because every germ can be represented on a singleton where the group is zero. Hence all stalks are zero although $\mathcal P(X)$ is not.
\end{solution}

\paragraph{What happened mathematically.}
The sheaf locality axiom is exactly what prevents global ``ghost sections'' invisible to every stalk.

\paragraph{Extensions.}
Relate this phenomenon to the canonical map from a presheaf to its sheafification.

\paragraph{Provenance.}
Legacy-derived from the stalk-detection proposition in the preferred sheaf notes.
\else
\paragraph{Provenance.}
Legacy-derived from the stalk-detection proposition in the preferred sheaf notes.
\fi
